% Combined UK Patent Application — Specification
% Application 2: BSDT + UDL
% Applicant: Odeyemi Olusegun Israel
% Date: 4 March 2026

\documentclass[12pt,a4paper]{article}
\usepackage[margin=2.5cm]{geometry}
\usepackage{amsmath,amssymb}
\usepackage{graphicx}
\usepackage{enumitem}
\usepackage{hyperref}
\usepackage{titlesec}
\usepackage{fancyhdr}
\usepackage{booktabs}
\setlength{\headheight}{14.5pt}

\pagestyle{fancy}
\fancyhf{}
\rhead{Patent Application}
\lhead{UDL/BSDT --- Odeyemi O.I.}
\rfoot{Page \thepage}

\titleformat{\section}{\large\bfseries}{\thesection.}{0.5em}{}
\titleformat{\subsection}{\normalsize\bfseries}{\thesubsection.}{0.5em}{}
\titleformat{\subsubsection}{\normalsize\itshape}{\thesubsubsection.}{0.5em}{}

\begin{document}

\begin{center}
{\LARGE\bfseries PATENT SPECIFICATION}\\[1em]
{\Large Multi-Spectrum Anomaly Detection and Failure-Mode\\
Correction System for Machine Learning Classifiers}\\[2em]
{\large Applicant: Odeyemi Olusegun Israel}\\[0.5em]
{\large Inventor: Odeyemi Olusegun Israel}\\[0.5em]
{\large Date of Filing: 4 March 2026}\\[0.5em]
{\large Application Number: [TO BE ASSIGNED]}\\[0.3em]
{\normalsize \textit{Filed pursuant to the Patents Act 1977}}
\end{center}

\vspace{2em}

% ============================================================
\section{TECHNICAL FIELD}
% ============================================================

The present invention relates to unsupervised anomaly detection and
post-hoc correction of deployed machine learning classifiers. More
particularly, it provides a two-stage framework in which a Universal
Deviation Layer (UDL) simultaneously measures deviations across multiple
independent mathematical law domains to construct an attribution-rich
anomaly tensor, and a Blind Spot Decomposition Theory (BSDT) layer uses
that tensor to detect and correct structural failure modes in any deployed
classification model without modifying the model's parameters. Applications
include fraud detection in financial, aerospace, and medical domains.

% ============================================================
\section{BACKGROUND}
% ============================================================

\subsection{Technical Problems}

\subsubsection{Limitation of Single-Perspective Anomaly Detection}

Standard anomaly detectors (Isolation Forest, Local Outlier Factor,
One-Class SVM, autoencoders) each produce a single scalar anomaly score.
They identify \emph{that} an observation is anomalous but not
\emph{which mathematical aspect} makes it anomalous. Because each detector
imposes its own implicit representation of normality, an observation
appearing normal under one representation may simultaneously violate a
different mathematical law. No single detector spans the full space of
mathematical anomaly, and none provides per-dimension attribution.
Furthermore, achieving high AUC (area under the ROC curve) does not
guarantee high coverage---anomaly scores may concentrate in a narrow band,
failing to discriminate among subtypes (the ``coverage--AUC gap'').

\subsubsection{Structural Failure Modes in Deployed Classifiers}

Supervised machine learning models deployed for fraud detection fail to
detect certain fraudulent observations for structural reasons arising from
specific, identifiable deficiencies in the model's representation of data
space. Retraining addresses failures only if fresh labelled data is
available, which it typically is not (fewer than 0.1\% of blockchain
transactions are labelled as illicit). Ensemble methods do not diagnose
\emph{why} specific instances are missed and may share the same structural
blind spot. Cost-sensitive reweighting is a global adjustment that does not
distinguish between failure modes. Classical bias-variance decomposition
operates in error space, not in feature space, providing no per-instance
diagnostic information of why a specific observation was missed.

\subsection{Prior Art Limitations}

Known prior art includes:
\begin{itemize}
\item Isolation Forest (Liu et al., 2008); Local Outlier Factor (Breunig
et al., 2000); One-Class SVM (Sch\"{o}lkopf et al., 1999): single-score
general-purpose detectors providing no attribution.
\item Deep autoencoders and out-of-distribution detection methods (ODIN,
energy-based OOD): based on a single latent representation; no multi-law
decomposition.
\item SHAP and LIME: explain individual predictions in terms of input
features but do not identify structural failure patterns or provide
correction mechanisms.
\item Conformal prediction: provides calibrated prediction sets but does not
diagnose failure modes.
\end{itemize}

None provides simultaneous multi-law-domain anomaly tensor construction,
MDN decomposition, boundary-centred dimension magnification, four-component
failure mode decomposition, or post-hoc correction operators that operate
on a frozen model.

% ============================================================
\section{SUMMARY OF THE INVENTION}
% ============================================================

The invention provides a unified two-stage framework. The shared inventive
concept is \emph{projecting observations through mathematically independent
law domains to produce an attribution-rich deviation decomposition, and
using that decomposition to discriminate and correct anomalies and
classification failures without model retraining}. The primary technical
contributions are:

\begin{enumerate}[label=(\roman*)]
\item The \textbf{Universal Deviation Principle}: a formal theorem that if
a family of spectrum operators $\{\Phi_k\}$ is deviation-separating (the
stacked Jacobian has full column rank), no anomalous observation can
simultaneously match normal data across all operators. This provides a
coverage guarantee absent from all single-perspective detectors.

\item A \textbf{multi-spectrum anomaly tensor} $T_{i,k,d}$ constructed by
projecting each observation through at least five independent mathematical
law domains: statistical, dynamical (chaos), spectral (frequency),
geometric, and exponential.

\item A \textbf{Magnitude--Direction--Novelty (MDN) decomposition}
separating how much ($r_i$), what type ($\hat{D}_i$), and how novel
($\theta_i$) each deviation is.

\item A \textbf{Boundary-Centred Dimension Magnifier} rescaling each
tensor dimension by its empirical class-separability, amplifying
well-separated dimensions and suppressing overlapping ones, using a
single global hyperparameter set.

\item A \textbf{four-component failure mode decomposition} (BSDT) of
the anomaly tensor into near-orthogonal components: Camouflage ($C$),
Feature Gap ($G$), Activity Anomaly ($A$), and Temporal Novelty ($T$),
serving as a per-instance diagnosis of why a deployed model fails.

\item \textbf{Self-calibrating weight computation} producing a composite
blind-spot score (MFLS) from the four components using mutual information,
Fisher variance ratio, Bayesian online updating, or gradient optimisation.

\item \textbf{Post-hoc correction operators} (additive, SignedLR,
QuadSurf, ExpGate) that apply targeted adjustments to the deployed model's
predicted probabilities based on the failure profile, recovering missed
detections without model retraining.

\item A formal \textbf{orthogonality verification} of the four components
(mean NMI $< 0.065$, all VIF $< 1.45$) confirming that each captures an
independent aspect of model failure.
\end{enumerate}

% ============================================================
\section{DETAILED DESCRIPTION}
% ============================================================

\subsection{Stage 1: Universal Deviation Layer (UDL)}
\label{sec:udl}

\subsubsection{Universal Deviation Principle}

\textbf{Definition (Deviation-Separating Family):} An operator family
$\{\Phi_k\}_{k=1}^{K}$ is \emph{deviation-separating} if the stacked
Jacobian:
\begin{equation}
J(x) = \begin{bmatrix}D\Phi_1(x)\\\vdots\\D\Phi_K(x)\end{bmatrix}
\end{equation}
has full column rank for all $x$ in the data domain.

\textbf{Theorem~1 (Local Deviation Guarantee):} If $\{\Phi_k\}$ is
deviation-separating, then for any anomalous point $x^*$ not identical
to the normal centroid, there exists at least one operator under which
$x^*$ deviates from the projection of normal data.

\textbf{Theorem~2 (Global Deviation Guarantee):} On a compact anomaly
set, a quantitative lower bound on minimum deviation across all operators
can be established.

\textbf{Theorem~3 (Finite-Sample Guarantee):} Given $n$ normal training
samples the detection guarantee holds with probability $\geq 1-\delta$
for a sample complexity depending on the operator family's covering number.

\textbf{Theorem~4 (Information-Optimal Fusion):} Optimal weights for
combining operator outputs are the solution to a semidefinite programme.

\textbf{Theorem~5 (Greedy Operator Selection):} A greedy algorithm
selecting the most informative operator subset achieves a $(1-1/e)$
approximation to the optimal subset under a submodular information
criterion.

\subsubsection{Spectrum Operators}

At least five mathematical law domains are instantiated as spectrum
operators $\Phi_k$:

\paragraph{Statistical Operator ($\Phi_{\mathrm{stat}}$).}
Measures distributional deviations: Shannon entropy of the local
neighbourhood; Kullback--Leibler divergence from the reference normal
distribution; Hellinger distance; skewness and kurtosis of local feature
distributions.

\paragraph{Dynamical Operator ($\Phi_{\mathrm{chaos}}$).}
Measures temporal dynamics: Lyapunov exponent from the temporal context;
recurrence rate from recurrence plot analysis; approximate entropy
measuring regularity.

\paragraph{Spectral Operator ($\Phi_{\mathrm{freq}}$).}
Measures frequency-domain deviations: power spectral density divergence;
frequency ratio (low-frequency to high-frequency energy); spectral entropy;
spectral centroid shift.

\paragraph{Geometric Operator ($\Phi_{\mathrm{geom}}$).}
Measures geometric relationships: Mahalanobis distance from the reference
centroid; cosine similarity to the reference direction; norm ratio.

\paragraph{Exponential Operator ($\Phi_{\mathrm{exp}}$).}
Amplifies small deviations invisible under linear representations:
\begin{equation}
\Phi_{\mathrm{exp}}(x) = \exp\!\left(\alpha\cdot\frac{x-\mu}{\sigma}
\right)
\end{equation}
where $\mu$,$\sigma$ are reference statistics.

In an enhanced embodiment, up to 14 composable operators are available
including phase-space, topological, kernel, rank, and wavelet operators.
An operator correlation audit confirms that at most one pair exceeds the
redundancy threshold ($\rho<0.85$).

\subsubsection{Anomaly Tensor Construction}

Each observation $x_i$ is projected through all $K$ operators to yield
a full anomaly tensor:
\begin{equation}
T_{i,k,d} \quad\text{for observation }i,\ \text{operator }k,
\ \text{dimension }d
\end{equation}

\subsubsection{Magnitude--Direction--Novelty Decomposition}

The tensor is decomposed into three interpretable components:
\begin{enumerate}[label=(\alph*)]
\item \textbf{Magnitude} $r_i = \|T_i\|$: total deviation amount.
\item \textbf{Direction} $\hat{D}_i = T_i/\|T_i\|$: unit tensor
indicating deviation type.
\item \textbf{Novelty} $\theta_i$: angular distance of $\hat{D}_i$ from
reference deviation directions in training data.
\end{enumerate}

\subsubsection{Boundary-Centred Dimension Magnifier}

Each tensor dimension is rescaled:
\begin{equation}
r'_d = \tanh\!\left(k_d\cdot\frac{r_d - b_d}{\sigma_d}\right)
\end{equation}
where $b_d$ is the decision boundary in dimension $d$, $\sigma_d$ is the
boundary-region scale, and $k_d = \gamma\cdot\delta_d$ is the steepness
with $\delta_d$ computed from three complementary measures:
\begin{enumerate}
\item Cohen's $d$ (effect size) between normal and anomalous distributions;
\item class-distribution overlap area; and
\item per-dimension AUC.
\end{enumerate}

The magnifier amplifies dimensions where classes are well-separated and
suppresses dimensions where they overlap. In the preferred embodiment
$\gamma{=}5.0$ and QDA regularisation $\lambda{=}10^{-4}$ constitute
the single global hyperparameter set used across all datasets.

\subsubsection{Classification Stage}

Four classification strategies, with hybrid auto-selection:
\begin{enumerate}[label=(\alph*)]
\item \textbf{Fisher projection:} LDA projection followed by thresholding.
\item \textbf{QDA:} Quadratic discriminant analysis with regularised
per-class covariance.
\item \textbf{QDA-Magnified:} QDA on the magnifier-rescaled tensor
(preferred strategy).
\item \textbf{RankFusion:} Fully unsupervised rank-statistic fusion.
\end{enumerate}

Hybrid auto-selection evaluates all strategies via cross-validation and
selects the one maximising coverage while maintaining AUC above a floor.

\subsubsection{Centroid Estimation}

Seven centroid strategies (sample mean, trimmed mean, geometric median,
MCD, medoid, Huber M-estimator, density-weighted centroid) are provided,
with auto-selection minimising within-class scatter.

\subsubsection{Coverage Metric}

\begin{equation}
\mathrm{Coverage@}k
= \frac{|\{x : \mathrm{rank}(x)\leq k,\ x\in\mathrm{anomaly}\}|}
{|\mathrm{anomaly}|}
\end{equation}

Coverage addresses the AUC--coverage gap. In the preferred embodiment
mean coverage exceeds 91\%.

\subsection{Stage 2: Blind Spot Decomposition Theory (BSDT) Layer}
\label{sec:bsdt}

The BSDT layer receives the anomaly tensor from Stage~1 (or may operate
directly on classifier outputs) and decomposes each observation's failure
profile into four near-orthogonal components.

\subsubsection{Relationship to Stage 1}

The BSDT four components correspond to specialised instantiations of the
UDL operator domains applied to a deployed classifier's feature space.
BSDT Camouflage ($C$) corresponds to the geometric operator applied to
the legitimate-class centroid; Feature Gap ($G$) corresponds to the
statistical operator measuring informational sparsity; Activity Anomaly
($A$) corresponds to the dynamical operator measuring deviation from known
fraud activity profiles; Temporal Novelty ($T$) corresponds to the spectral
and dynamical operators applied over the temporal axis of the training
distribution.

\subsubsection{Camouflage Component ($C$)}

Measures proximity of the observation to the legitimate-class distribution:
\begin{equation}
C(x) = 1 - \frac{\|x-\mu_{\mathrm{legit}}\|_2}{d_{\max}}
\end{equation}
where $\mu_{\mathrm{legit}}$ is the legitimate-class centroid and
$d_{\max}$ is the maximum centroid distance in the training data.
$C(x)\approx 1$ indicates high camouflage: the observation mimics
legitimate patterns, causing the model to under-estimate its suspicion.

\subsubsection{Feature Gap Component ($G$)}

Measures informational sparsity:
\begin{equation}
G(x) = \frac{|\{j : |x_j|<\varepsilon\}|}{d}
\end{equation}
where $d$ is the feature dimension and $\varepsilon=10^{-8}$.
High $G(x)$ indicates that many features are uninformative, depriving the
model of discriminative signal.

\subsubsection{Activity Anomaly Component ($A$)}

Measures deviation of transactional activity from the activity profile of
previously detected fraud:
\begin{equation}
A(x) = \sigma\!\left(
\frac{\log(1+|\mathrm{activity}(x)|)-\mu_{\mathrm{ref}}}
{\sigma_{\mathrm{ref}}}\right)
\end{equation}
where $\mathrm{activity}(x)$ is a transactional activity metric,
$\mu_{\mathrm{ref}}$ and $\sigma_{\mathrm{ref}}$ are statistics of the
reference detected-anomaly population, and $\sigma$ is the logistic
sigmoid. This component captures fraud operating at activity scales the
model has not learned to associate with fraud.

\subsubsection{Temporal Novelty Component ($T$)}

Measures the degree to which the observation is temporally novel relative
to training data:
\begin{equation}
T(x) = \sigma\!\left(\tfrac{1}{2}(\bar{m}(x)-2)\right)
\end{equation}
where $\bar{m}(x)$ is the mean distance to the $k$ nearest temporal
neighbours in the training set, normalised by the temporal span.

\subsubsection{Self-Calibrating Weight Computation}

The four components are aggregated into a composite blind-spot score (MFLS):
\begin{equation}
\mathrm{MFLS}(x) = \sum_{i\in\{C,G,A,T\}} w_i\cdot S_i(x)
\end{equation}

The weights $w_i$ (normalised to sum to 1) are computed automatically by
one or more of:
\begin{enumerate}[label=(\alph*)]
\item Mutual Information between each component and the binary label on a
held-out calibration set;
\item Fisher Variance Ratio (between-class to within-class variance);
\item Bayesian Online updating as labelled observations arrive; or
\item Gradient-based optimisation maximising detection performance.
\end{enumerate}

\subsubsection{Correction Operators}

Four correction operators apply targeted adjustments to the frozen model's
predicted probabilities $p(x)$:

\paragraph{Additive Correction.}
\begin{equation}
p^*(x) = p(x) + \lambda\cdot\mathrm{MFLS}(x)\cdot(1-p(x))\cdot
\mathbb{1}[p(x)<\tau]
\end{equation}

\paragraph{Logistic Regression (SignedLR).}
\begin{equation}
p^*(x) = \mathrm{clip}\!\left(p(x)+\beta_0+\textstyle\sum_i\beta_i S_i(x),
\ 0,\ 1\right)
\end{equation}

\paragraph{Quadratic Surface (QuadSurf).}
\begin{equation}
p^*(x) = \mathrm{clip}\!\left(p(x)+\phi_2([C,G,A,T])^\top\hat{\beta},
\ 0,\ 1\right)
\end{equation}
where $\phi_2$ is the degree-2 polynomial feature map (raw components,
squared components, pairwise products) and $\hat{\beta}$ are coefficients
fitted by ridge regression ($\alpha{=}1.0$).

\paragraph{Exponential Gate (ExpGate).}
\begin{equation}
p^*(x) = \mathrm{clip}\!\left(p(x)+g(x)\cdot q(x),\ 0,\ 1\right)
\end{equation}
where $q(x)$ is the QuadSurf output and
$g(x)=\sigma(v^\top[C,G,A,T]+b)$ is a learned sigmoid gate.

\subsubsection{Seven Supplementary Features}

In an enhanced embodiment, seven supplementary features supplement the
core components for use by correction operators: the four raw components
$C,G,A,T$; the composite MFLS score; the predicted-vs-actual residual
$|p(x)-y|$ (available in supervised calibration); and the interaction
term $C(x)\cdot A(x)$.

\subsubsection{Orthogonality Verification}

The four components are verified near-orthogonal by three independent
measures:
\begin{itemize}
\item Normalised Mutual Information: mean NMI $\in[0.051,0.064]$ across all
pairs, well below the 0.30 redundancy threshold.
\item Variance Inflation Factor: all VIF $<1.45$, confirming negligible
multicollinearity.
\item PCA: all four components carry meaningful variance; none is dominated.
\end{itemize}

\subsubsection{Cross-Domain Transfer}

The adaptive damping function derived from the composite blind-spot energy:
\begin{equation}
\gamma^*(E)=\frac{E}{E+\theta}
\end{equation}
is also applicable as a viscosity regularisation term in Navier--Stokes
fluid dynamics simulations, representing a novel finance-to-physics transfer
and linking the BSDT framework to the systemic risk controller described in
related Application~1 (AMTTP/AF).

\subsection{Cross-Domain Performance}

Domain-dependent dominant components confirm genuine structural differences
across application domains:

\begin{center}
\begin{tabular}{lll}
\toprule
\textbf{Domain} & \textbf{Dataset} & \textbf{Dominant Component} \\
\midrule
Blockchain & Elliptic & Temporal Novelty ($T$) \\
Blockchain & Ethereum Addresses & Temporal Novelty ($T$) \\
Banking & Credit Card Fraud & Activity Anomaly ($A$) \\
Banking & IEEE-CIS & Activity Anomaly ($A$) \\
Aerospace & Shuttle & Camouflage ($C$) \\
Medical & Mammography & Camouflage ($C$) \\
Pattern Recognition & Pendigits & Feature Gap ($G$) \\
\bottomrule
\end{tabular}
\end{center}

\subsection{Performance Characteristics}

\textbf{UDL --- Single hyperparameter set ($\gamma{=}5.0$,
$\lambda{=}10^{-4}$):}

\begin{center}
\begin{tabular}{lcc}
\toprule
\textbf{Dataset} & \textbf{QDA-Magnified AUC} & \textbf{Hybrid AUC} \\
\midrule
Synthetic & 1.000 & 0.999 \\
Mimic (camouflage adversarial) & 1.000 & 0.999 \\
Mammography & 0.947 & 0.911 \\
Shuttle & 0.989 & 0.991 \\
Pendigits & 0.777 & 0.960 \\
\textbf{Mean} & \textbf{0.943} & \textbf{0.972} \\
\bottomrule
\end{tabular}
\end{center}

\textbf{BSDT --- Evaluated across 792 variant--model--dataset
combinations:}

\begin{center}
\begin{tabular}{ll}
\toprule
\textbf{Metric} & \textbf{Value} \\
\midrule
Best variant (QuadSurf) mean F1 & 0.787 \quad ($+43\%$ relative) \\
False positive reduction & 88.1\% (44,740 $\to$ 5,345) \\
Maximum recall recovery & $+84.4$ percentage points \\
Shuttle rescue (F1) & $0.134\to 0.968$ \\
Shuttle FP reduction & 99.94\% \\
Zero-missed-fraud configurations & 10\,/\,12 model--dataset combos \\
Cross-domain credit card MFLS AUC & 0.885 \\
Cross-domain aerospace MFLS AUC & 0.982 \\
Cross-domain medical imaging MFLS AUC & 0.916 \\
\bottomrule
\end{tabular}
\end{center}

% ============================================================
\section{CLAIMS}
% ============================================================

\noindent\textit{The following claims define the scope of protection
sought. Where a claim refers to another claim, the dependency is by claim
number. Features of any dependent claim may be combined with any independent
claim or with other dependent claims.}
\vspace{0.8em}

\begin{enumerate}[label=\textbf{\arabic*.}]

% --- Independent Claim 1: UDL Detection Method ---
\item A computer-implemented method for detecting and characterising anomalies
in data, the method comprising:
\begin{enumerate}[label=(\alph*)]
\item receiving an input observation comprising a feature vector;

\item projecting the observation through a plurality of spectrum operators,
each computing deviation measurements under a different mathematical law
domain, the plurality comprising at least:
\begin{enumerate}[label=(\roman*)]
\item a statistical operator measuring distributional deviations;
\item a dynamical operator measuring temporal dynamical deviations;
\item a spectral operator measuring frequency-domain deviations;
\item a geometric operator measuring geometric relationship deviations; and
\item an exponential operator measuring deviations under exponential
amplification;
\end{enumerate}

\item constructing a multi-dimensional anomaly tensor from the deviation
measurements across all spectrum operators;

\item decomposing the anomaly tensor into magnitude, direction, and novelty
components, wherein the magnitude measures total deviation amount, the
direction measures deviation type, and the novelty measures angular distance
of the deviation direction from reference directions in the training data;

\item applying a boundary-centred dimension magnifier that rescales each
dimension of the anomaly tensor by a per-dimension steepness coefficient
$k_d = \gamma \cdot \delta_d$ where $\delta_d$ is a discriminability score
computed from empirical class-separability in that dimension; and

\item classifying the rescaled anomaly tensor to produce an anomaly score
and per-law-domain attribution identifying which mathematical law domains
contribute to the anomaly determination.
\end{enumerate}

% --- Independent Claim 2: Combined UDL + BSDT System ---
\item A computer-implemented system for detecting anomalies and correcting
failure modes in deployed machine learning classifiers, the system
comprising:
\begin{enumerate}[label=(\alph*)]
\item a plurality of spectrum operator modules collectively satisfying a
deviation-separating condition such that the stacked Jacobian of the
operators has full column rank, each module computing deviation measurements
under a different mathematical law domain;

\item a tensor construction module configured to assemble the deviation
measurements into a multi-dimensional anomaly tensor;

\item a decomposition module configured to decompose the anomaly tensor into
magnitude, direction, and novelty components;

\item a magnifier module configured to rescale each dimension of the anomaly
tensor using a boundary-centred rescaling function with per-dimension
steepness computed from empirical class separability;

\item a classifier module configured to produce an anomaly score and
per-law-domain attribution from the rescaled tensor;

\item a failure-mode decomposition module configured to compute, for each
input observation and an associated predicted probability output by a
deployed frozen classification model, four near-orthogonal failure mode
components comprising:
\begin{enumerate}[label=(\roman*)]
\item a camouflage component measuring proximity of the observation to a
legitimate-class centroid;
\item a feature gap component measuring the proportion of near-zero features;
\item an activity anomaly component measuring deviation of a transactional
activity metric from a reference detected-anomaly population; and
\item a temporal novelty component measuring temporal distance from training
data patterns;
\end{enumerate}

\item a weighting module configured to compute component weights using one
or more of mutual information, Fisher variance ratio, Bayesian online
updating, or gradient-based optimisation, and to aggregate the weighted
components into a composite blind-spot score; and

\item a correction module configured to apply one or more correction
operators to the predicted probabilities of the deployed model based on
the computed component values, to produce corrected probabilities;
\end{enumerate}
wherein the failure-mode decomposition module and correction module operate
on a frozen deployed model without requiring modification or retraining of
that model.

% --- Independent Claim 3: BSDT Correction Method ---
\item A computer-implemented method for detecting and correcting structural
failure modes in a deployed frozen machine learning classification model
without modifying the model, the method comprising:
\begin{enumerate}[label=(\alph*)]
\item receiving a feature vector $x$ and a predicted probability $p(x)$
output by the deployed model;

\item computing a camouflage component $C(x)$ measuring proximity of the
feature vector to a legitimate-class centroid in feature space;

\item computing a feature gap component $G(x)$ measuring the proportion of
features having values below a near-zero threshold;

\item computing an activity anomaly component $A(x)$ measuring deviation
of a transactional activity metric from the activity profile of a reference
detected-anomaly population;

\item computing a temporal novelty component $T(x)$ measuring the degree
to which the observation is temporally distant from patterns present in
the training data;

\item computing a composite blind-spot score by aggregating the four
components using automatically calibrated weights; and

\item applying a correction operator to the predicted probability $p(x)$
based on values of one or more computed components to produce a corrected
probability $p^*(x)$.
\end{enumerate}

% --- Dependent Claims ---
\item A method according to claim 1, wherein the boundary-centred dimension
magnifier applies the rescaling function:
\[
r'_d = \tanh\!\left(k_d\cdot\frac{r_d - b_d}{\sigma_d}\right)
\]
where $k_d=\gamma\cdot\delta_d$, $\gamma$ is a global steepness parameter,
and $\delta_d$ is a per-dimension discriminability score computed from
Cohen's $d$, class-distribution overlap, and per-dimension AUC.

\item A method according to claim 1, wherein classifying the rescaled tensor
uses quadratic discriminant analysis with regularised per-class covariance
matrices applied to the rescaled anomaly tensor.

\item A method according to claim 1, further comprising a hybrid
auto-selection step that evaluates a plurality of classification strategies
on cross-validation folds and selects the strategy maximising a coverage
metric while maintaining AUC above a minimum threshold, wherein the coverage
metric is:
\[
\mathrm{Coverage@}k = \frac{|\{x : \mathrm{rank}(x)\leq k,\
x\in\mathrm{anomaly}\}|}{|\mathrm{anomaly}|}
\]

\item A method or system according to any of claims 1 to 6, further
comprising an operator correlation audit verifying independence of the
spectrum operators by confirming pairwise correlation coefficients are
below a redundancy threshold.

\item A method or system according to any of claims 1 to 7, further
comprising a centroid auto-selection module selecting from a plurality of
centroid estimation strategies the strategy minimising within-class scatter
of normal data in the rescaled space.

\item A method or system according to any of claims 1 to 8, wherein a
single global hyperparameter set is used across all datasets without
per-dataset tuning.

\item A method or system according to any of claims 1 to 9, further
comprising a greedy operator selection step that selects a subset of
operators from a larger candidate set achieving a $(1-1/e)$ approximation
to the optimal subset under a submodular information criterion.

\item A method according to claim 3, wherein the camouflage component is:
\[
C(x) = 1 - \frac{\|x-\mu_{\mathrm{legit}}\|_2}{d_{\max}}
\]
\item A method according to claim 3, wherein the feature gap component is:
\[
G(x) = \frac{|\{j : |x_j|<\varepsilon\}|}{d}
\]
\item A method according to claim 3, wherein the activity anomaly component is:
\[
A(x) = \sigma\!\left(\frac{\log(1+|\mathrm{activity}(x)|)-\mu_{\mathrm{ref}}}
{\sigma_{\mathrm{ref}}}\right)
\]

\item A method according to claim 3, wherein the correction operator is a
QuadSurf correction:
\[
p^*(x) = \mathrm{clip}\!\left(p(x)+\phi_2([C,G,A,T])^\top\hat{\beta},
\ 0,\ 1\right)
\]
where $\phi_2$ is the degree-2 polynomial feature map and $\hat{\beta}$
are ridge-regression coefficients.

\item A method according to claim 3, wherein the correction operator is an
additive correction:
\[
p^*(x) = p(x) + \lambda\cdot\mathrm{MFLS}(x)\cdot(1-p(x))
\cdot\mathbb{1}[p(x)<\tau]
\]

\item A method according to claim 3, wherein the correction operator is an
exponential gating correction comprising a QuadSurf output modulated by a
learned sigmoid gate $g(x)=\sigma(v^\top[C,G,A,T]+b)$.

\item A method or system according to any of claims 3 to 16, wherein the
four components are verified near-orthogonal by having normalised mutual
information below 0.30 for all pairs and variance inflation factors below
2.0 for all components.

\item A method or system according to any of claims 3 to 17, wherein the
dominant component varies by application domain: temporal novelty in
blockchain domains, activity anomaly in banking domains, and camouflage
in aerospace and medical domains.

\item A method or system according to any of claims 3 to 18, further
comprising the step of computing seven supplementary features comprising the
four raw components, the composite blind-spot score, a predicted-vs-actual
residual, and a camouflage-activity interaction term.

\item A method or system according to any preceding claim, wherein the
adaptive damping function $\gamma^*(E)=E/(E+\theta)$ derived from the
composite blind-spot energy is applied as a viscosity regularisation term
in a computational fluid dynamics simulation.

\end{enumerate}

% ============================================================
\section{DRAWINGS}
% ============================================================

\textit{[To be filed separately:
(1) Architecture diagram of the two-stage UDL + BSDT pipeline;
(2) Diagram of the five spectrum operators and the anomaly tensor
construction;
(3) Visualisation of the MDN magnitude--direction--novelty decomposition;
(4) Schematic of the boundary-centred dimension magnifier with per-dimension
discriminability;
(5) Operator correlation matrix confirming independence;
(6) Four-component BSDT decomposition diagram;
(7) Flow diagram of the correction operator pipeline;
(8) NMI matrix and VIF chart confirming near-orthogonality;
(9) Cross-domain performance comparison chart;
(10) Schematic of BSDT operating on a frozen deployed classifier.]}

% ============================================================
\newpage
\section*{ABSTRACT}
% ============================================================

A two-stage computer-implemented framework detects anomalies and corrects
structural failure modes in deployed machine learning classifiers without
model retraining. Stage~1 (Universal Deviation Layer) projects each
observation through at least five independent mathematical law domains,
constructs a multi-dimensional anomaly tensor, decomposes it into magnitude,
direction and novelty components, and applies a boundary-centred dimension
magnifier to produce an attributed anomaly score. A formal deviation
guarantee proves that no anomaly can simultaneously match normal data across
all law domains. Stage~2 (Blind Spot Decomposition Theory) decomposes each
observation's failure profile into four near-orthogonal components:
camouflage, feature gap, activity anomaly and temporal novelty. A composite
blind-spot score with self-calibrating weights drives post-hoc correction
operators that adjust a frozen model's predicted probabilities, recovering
up to 84.4 percentage points of recall. The framework is validated across
financial, aerospace and medical datasets using a single hyperparameter set.

\end{document}
